\section{Framework Design}
\label{sec:design}

CT-KAT (invoked as \texttt{ctkat}) is a fail-closed CT screening framework for
PQC. It binds three complementary checkers under one declarative configuration,
sweeps every harness across a compiler${}\times{}$optimization build matrix, and
routes every finding through a default-deny triage taxonomy backed by a cited
registry of accepted variable-time behavior. The goal is not to prove the
absence of leaks --- no dynamic tool can --- but to make screening honest: each
checker covers a leak class the others miss, the verdict is reported per build,
and a detected CT failure is never silently promoted to a key-recovery leak. The
pipeline (Fig.~\ref{fig:pipeline}) spans configuration (\S\ref{ssec:config}),
checkers (\S\ref{ssec:checkers}), the build sweep (\S\ref{ssec:matrix}), and
triage (\S\ref{ssec:triage}).

\subsection{Unified Configuration and Pipeline}
\label{ssec:config}

One YAML file describes a run: the target harnesses (each a small driver calling
one cryptographic operation, e.g.\ a decapsulation), the secret input regions to
track, the named build configurations, and per-layer options. As the single
source of truth it fixes the build flags, tainted regions, and timing parameters
together, so the measured configuration cannot drift from the assumed one. Each
harness is compiled once per build cell (\S\ref{ssec:matrix}), screened by the
three layers (\S\ref{ssec:checkers}), and reduced by triage
(\S\ref{ssec:triage}) to one \texttt{verdict\_class}.

\subsection{Three Checker Layers}
\label{ssec:checkers}

CT side channels in PQC share no single observable signature, so CT-KAT composes
three checkers, each for a leak class the others cannot see.

\paragraph{Layer 1: structural CT check (Valgrind/Memcheck).}
The classic dynamic-taint approach, after ctgrind~\cite{ctgrind} on
Memcheck~\cite{valgrind}: secret inputs are marked undefined, and Memcheck
reports when a secret-derived value reaches a control-flow branch or a memory
address --- the constructs behind the canonical cache- and branch-timing
leaks~\cite{kocher}. It checks program \emph{structure}, not instruction
\emph{latency}, so a secret-dependent variable-latency division passes it clean.

\paragraph{Layer 2: variable-latency instruction scan (asm-scan).}
\texttt{asm-scan} closes that gap: a warn-only static scan over the emitted
assembly, run across optimization levels, flagging \emph{candidate}
variable-latency instructions --- chiefly integer division and remainder, whose
amd64 latency depends on operands. The KyberSlash class~\cite{kyberslash} is
exactly such a secret-derived division that Memcheck passes clean. It stays
warn-only because the divisor may be public; the public-versus-secret split is
resolved in triage (\S\ref{ssec:triage}).

\paragraph{Layer 3: black-box statistical timing (dudect).}
A black-box test after dudect~\cite{dudect}: it runs the harness on two input
classes and applies a Welch $t$-test to the cycle counts; a large $|t|$ means
the running time depends on the input class, hence the secret. Needing no model
of why a leak occurs, it catches timing dependence with no structural branch and
no recognizable variable-latency instruction. Being measurement-based it is
environment-sensitive: under our virtualized image it can report
environment-induced noise, which it surfaces loudly. The three layers detect
disjoint leak classes --- structural, instruction latency, end-to-end timing ---
so their union is the coverage.

\subsection{Build-Configuration Sweep (ct-matrix)}
\label{ssec:matrix}

A constant-time property belongs to the \emph{compiled} program: the compiler
may remove a secret-dependent branch at one optimization level, or inline a
function so a finding lands in a different frame, making a single-build verdict
unsound. \texttt{ct-matrix} recompiles each harness under every named build
configuration and re-runs the \emph{same} structural check per cell, reporting
the verdict as a function of the build. The matrix is indexed by compiler and
named cflags; our image provides gcc~13.3.0 and clang~18.1.3, giving up to six
cells. The three gcc configurations are \texttt{debug}
(\texttt{-O0 -g -fno-inline -fno-omit-frame-pointer}, disabling inlining so
findings name their precise leak-site function), \texttt{release}
(\texttt{-O2 -g -fno-omit-frame-pointer -fno-lto}), and \texttt{size}
(\texttt{-Os}, otherwise as release). Two phenomena make the sweep load-bearing:
a branch present at \texttt{-O0} can be removed at \texttt{-O2}/\texttt{-Os},
flipping the verdict; and inlining blurs attribution, so \texttt{-O0 -fno-inline}
names the true leak-site function while optimized builds merge it into a parent
frame. CT-KAT unions the leak-site functions and treats the debug build as
authoritative for triage while reporting build-dependence.

\subsection{Default-Deny Triage Taxonomy}
\label{ssec:triage}

Over-claiming every structural failure as a key-recovery leak is dangerous for
PQC, since several schemes are variable-time \emph{by design} --- e.g.\ the
rejection-sampling loops of ML-DSA~\cite{fips204,dilithium}, whose iteration
count depends on a fresh per-signature nonce. CT-KAT assigns each harness exactly
one \texttt{verdict\_class}:
\begin{description}
  \item[\texttt{robust}:] all checkers pass across the matrix.
  \item[\texttt{ct-clean-untriaged}:] structural pass; another layer's candidate
        (e.g.\ a division) is untriaged.
  \item[\texttt{varlat-secret-risk}:] structural pass, but \texttt{asm-scan}
        flags a \emph{secret-derived} variable-latency instruction (the
        KyberSlash case).
  \item[\texttt{build-sensitive-ct}:] the structural verdict is unstable across
        the matrix.
  \item[\texttt{accepted-variable-time}:] structural fail, but \emph{every}
        leak-site function is a registered, cited design behavior.
  \item[\texttt{needs-analysis}:] structural fail with some unregistered
        leak-site function; a human must rule.
  \item[\texttt{ct-leak}:] a confirmed leak, reachable \emph{only} by explicit
        human verdict; never assigned automatically.
\end{description}

The \texttt{accepted-variable-time}/\texttt{needs-analysis} boundary is a cited
registry (\texttt{docs/accepted\_variable\_time.md}): each entry records a
function whose secret-dependent timing is an analyzed-safe design behavior, with
a checkable citation. If \emph{all} leak-site functions are registered the
harness is \texttt{accepted-variable-time}; if \emph{any} is unrecognized it
stays \texttt{needs-analysis} --- never silently accepted, never silently
promoted to a leak. Section~\ref{sec:results-coverage} shows this default-deny
invariant at work on real ML-DSA-65.
